\documentclass[manuscript,screen 12pt, nonacm]{acmart}
\let\Bbbk\relax % Fix for amssymb clash 
\usepackage{vmlmacros}
\AtBeginDocument{%
  \providecommand\BibTeX{{%
    \normalfont B\kern-0.5em{\scshape i\kern-0.25em b}\kern-0.8em\TeX}}}
\usepackage{outlines}
\setlength{\headheight}{14.0pt}
\setlength{\footskip}{13.3pt}
\title{An Alternative to Pattern Matching, Inspired by Verse}
\author{Roger Burtonpatel}
\email{roger.burtonpatel@tufts.edu}
\affiliation{%
  \institution{Tufts University}
  \streetaddress{419 Boston Ave}
  \city{Medford}
  \state{Massachusetts}
  \country{USA}
  \postcode{02155}
}
\begin{document}



\section{Pattern Matching and its Extensions}
\label{pmandextensions}

    Pattern matching lets programmers examine and deconstruct data by matching them
    against patterns. 

    % 1-2 more sentences. PM is beloved and well-established. 


    % PM Terms: maybe 

    \hyperref[fig:pmarea]{\tt{pm\_area}} provides an opportunity to introduce a
    few terms that are common in pattern matching.
    \hyperref[fig:pmarea]{\tt{pm\_area}}~is a classic example of where pattern
    matching most commonly occurs: within a~\it{case} expression. A~\it{case}
    expression tests a~\it{scrutinee} (\tt{sh}) against a list of~\it{branches}.
    Each branch contains a pattern on the left-hand side (\tt{SQUARE s}, etc.)
    and an expression on the right-hand side (\tt{s * s}, etc.). When a pattern
    matches the result of evaluating the scrutinee, the program
    evaluates the right-hand side of the respective branch.\footnote{OCaml,
    which you'll see in future sections, calls~\it{case}~\tt{match}. Some
    literature~\citep{guardproposal} calls this a~\it{head expression}. I~follow
    the example of~\citet{bpc} and~\citet{maranget} in calling the things
    \it{case} and~\it{scrutinee}. Any of these terms does the job.} 

\subsection{Popular extensions to pattern matching}
\label{extensions}

    Extensions to pattern matching simplify cases that are otherwise
    troublesome. 
    % No nice properties- they increase brevity and closeness to algebraic laws.
    % citation is good here. 
    
    % Specifically, extensions help restore Nice~Properties
    % \hyperref[p1]{Lawlike} and~\hyperref[p2]{Single copy} in cases where pattern
    % matching fails to satisfy them. 
    
    In this section, I~illustrate several such cases, and I~demonstrate how
    extensions help programmers write code that adheres to the Nice~Properties.
    The three extensions I~describe are those commonly found in the literature
    and implemented in compilers: side conditions, pattern guards, and
    or-patterns. 
    
    To denote pattern matching~\it{without} extensions, I~coin the term~\it{bare
    pattern matching}. In bare pattern matching, a pattern has one of two
    syntactic forms: a name or an application (of a value constructor to zero or
    more patterns). 


% Subject: Extensions to pattern matching. 

\subsubsection{Side conditions}

    First, I~illustrate why programmers want~\it{side conditions}, an extension
    to pattern matching common in most popular functional programming languages,
    including OCaml, Erlang, Scala, and~Haskell\footnote{I~use the term~\it{side
    conditions} to refer to a pattern followed by a Boolean expression. Some
    languages call this a~\it{guard}, which I~use to describe a different, more
    powerful extension to pattern matching in Section~\ref{guards}. Haskell has
    \it{only} guards, but a Boolean guard~\underline{is} a side condition.}. 
    
    I~define a (rather silly) function~\tt{exclaimTall} in OCaml on
    \tt{standard\_shape}s. I~have to translate our~\tt{standard\_shape} datatype to OCaml, and
    while I'm at it, I~write the type and algebraic laws for
    \tt{exclaimTall}:

    \begin{minipage}[t]{\textwidth}        
        \centering 
        \begin{verbatim}

type standard_shape = Square of float 
                    | Triangle of float * float
                    | Trapezoid of float * float * float

exclaimTall : standard_shape -> string 

exclaimTall (Square s)              == "Wow! That's a sizeable square!", 
                                        where s > 100.0
exclaimTall (Triangle (w, h))       == "Goodness! Towering triangle!",
                                        where h > 100.0
exclaimTall (Trapezoid (b1, b2, h)) == "Zounds! Tremendous trapezoid!", 
                                        where h > 100.0
exclaimTall sh                      == "Your shape is small.", 
                                        otherwise
    \end{verbatim}
    \end{minipage}

    Armed with pattern matching, I~implement~\tt{exclaimTall} in OCaml
    (Figure~\ref{fig:ifexclaimtall}).

    \begin{figure}[ht]
        \begin{verbatim}
            let exclaimTall sh =
            match sh with 
              | Square s -> if s > 100.0 
                            then "Wow! That's a sizeable square!"
                            else "Your shape is small." 
              | Triangle (_, h) -> 
                            if h > 100.0 
                            then "Goodness! Towering triangle!"
                            else "Your shape is small." 
              | Trapezoid (_, _, h) -> 
                            if h > 100.0
                            then "Zounds! Tremendous trapezoid!"
                            else "Your shape is small." 
            \end{verbatim}    
        \caption{An invented function~\tt{exclaimTall} in OCaml combines pattern
        matching with an~\tt{if} expression, and is not very pretty.}   
        \Description{An implementation of a function exclaimTall in OCaml, which
        uses pattern matching and an if statement.}
        \label{fig:ifexclaimtall}
    \end{figure}
    
    Here, I'm using the special variable~\tt{\_}---that's the underscore
    character, a wildcard pattern---to indicate that I~don't care about the
    bindings of a pattern. 

    I~am not thrilled with this code. It gets the job done, but it fails to
    adhere to the Nice~Properties of~\hyperref[p1]{Lawlike}
    and~\hyperref[p2]{Single copy}: the code does not look like the algebraic
    laws, and it duplicates right-hand side,~\tt{"Your shape is small"}, three
    times. I~find the code unpleasant to read, too: the actual “good” return
    values of the function, the exclamatory strings, are gummed up in the middle
    of the~\tt{if-then-else} expressions.
    
    Fortunately, this code can be simplified by using the~\tt{standard\_shape} patterns
    with a~\it{side condition}, i.e., a syntactic form for “match a pattern
    \it{and} a Boolean condition.” The~\tt{when} keyword in OCaml provides such
    a form, as seen in Figure~\ref{fig:whenexclaimtall}.
        
        \begin{figure}[]
            \begin{verbatim}
                let exclaimTall sh =
                  match sh with 
                    | Square s when s > 100.0 ->
                            "Wow! That's a sizeable square!"
                    | Triangle (_, h) when h > 100.0 ->
                            "Goodness! Towering triangle!"
                    | Trapezoid (_, _, h) when h > 100.0 -> 
                            "Zounds! Tremendous trapezoid!"
                    | _ ->  "Your shape is small." 
                \end{verbatim}
            \caption{With a side condition,~\tt{exclaimTall} in OCaml becomes
            simpler and more adherent to the Nice~Properties.} 
            \Description{An implementation of a function exclaimTall in OCaml,
            which uses pattern matching and a side condition.}
            \label{fig:whenexclaimtall}
        \end{figure}

    A side condition streamlines pattern-and-Boolean cases and minimizes
    overhead, restoring~\hyperref[p1]{Lawlike} and~\hyperref[p2]{Single copy}.
    And a side condition can exploit bindings that emerge from the preceding
    pattern match. For instance, the~\tt{when} clauses in
    Figure~\ref{fig:whenexclaimtall} exploit names~\tt{s} and~\tt{h}, which are
    bound in the match of~\tt{sh} to~\tt{Square s},~\tt{Triangle (\_, h)}, and
    \tt{Trapezoid (\_,~\_, h)}, respectively. 

    Importantly, side conditions come at a cost: their inclusion means that a
    compiler enforcing the~\hyperref[p5]{Exhaustiveness} Nice~Property becomes
    an NP-hard problem, because it must now perform exhaustiveness analysis not
    only on patterns, but on arbitrary expressions. Modern compilers give a
    weaker form of exhaustiveness that only deals with patterns, and side
    conditions are worth the tradeoff for restoring the two most important of
    the Nice~Properties:~\hyperref[p1]{Lawlike} and~\hyperref[p2]{Single copy}.

    A side condition adds an extra “check”---in this case, a Boolean
    expression---to a pattern. But side conditions have a limitation: the check
    can make a decision based off of an expression evaluating to~\tt{true} or
    \tt{false}, but not an expression evaluating to, say, \tt{nil} or \tt{::}
    (cons). In the next section, I~use an example to~showcase when this
    limitation matters, and how another extension addresses it. 

    \subsubsection{Pattern guards}
    \label{guards}

    To highlight a common use of pattern guards to address such a limitation, I
    modify an example from~\citet{guardproposal}, the proposal for pattern
    guards in GHC. Suppose I~have an OCaml abstract data type of finite maps,
    with a lookup operation: 

    \begin{minipage}[t]{\textwidth}
        \centering 
        \begin{verbatim}
            lookup : finitemap -> int -> int option
        \end{verbatim}
    \end{minipage}
    Let's say I~want to perform three successive lookups, and call a “fail”
    function if~\it{any} of them returns~\tt{None}. Specifically, I~want a
    function with this type and algebraic laws: 

    \begin{minipage}[t]{\textwidth}
        \centering 
        \begin{verbatim}
          

            tripleLookup : finitemap -> int -> int

              tripleLookup rho x == z, where 
                                        lookup rho x == Some w
                                        lookup rho w == Some y
                                        lookup rho y == Some z
              tripleLookup rho x == handleFailure x, otherwise
            
              handleFailure : int -> int 

              handleFailure's implementation is unimportant.
                handleFailure (x : int) = ... some error-handling ... -> x  

        \end{verbatim}
    \end{minipage}

    To express this computation succinctly, the program needs to make decisions
    based on how successive computations match with patterns, but neither bare
    pattern matching nor side conditions give that flexibility. 
    
    Side conditions don't appear to help here, so I~try with bare pattern
    matching. Figure~\ref{fig:pmtriplelookup} shows how I~might implement
    \tt{tripleLookup} as such. 

    \begin{figure}[ht]
        \begin{verbatim}

          let tripleLookup (rho : finitemap) (x : int) =
              match lookup rho x with Some w -> 
                (match lookup rho w with Some y -> 
                  (match lookup rho y with Some z -> z
                  | _ -> handleFailure x)
                | _ -> handleFailure x)
              | _ -> handleFailure x
            \end{verbatim}
        \caption{\tt{tripleLookup} in OCaml with bare pattern matching breaks
         the Nice~Property of~\hyperref[p2]{Single copy}: avoiding duplicated
         code.} 
                    
        \Description{An implementation of the tripleLookup with three levels of
        nested pattern matching.}
        \label{fig:pmtriplelookup}
    \end{figure}

    Once again, the code works, but it's lost
    the~\hyperref[p1]{Lawlike}~and~\hyperref[p2]{Single copy} Nice~Properties by
    duplicating three calls to~\tt{handleFailure} and stuffing the screen full
    of syntax that distracts from the algebraic laws. Unfortunately, it's not
    obvious how a side condition could help us here, because we need pattern
    matching to extract and name internal values~\tt{w},~\tt{y}, and~\tt{z}
    from constructed data.

    To restore the Nice~Properties, I~introduce a more powerful extension to
    pattern matching:~\it{pattern guards}, a form of “smart pattern” in which
    intermediate patterns bind to expressions within a single branch of a
    \tt{case}. Pattern guards can make~\tt{tripleLookup} appear~\it{much}
    simpler, as shown in Figure~\ref{fig:guardtriplelookup}---which, since
    pattern guards aren't found in OCaml, is written in Haskell.

    \begin{figure}
        \begin{center}
        \begin{verbatim}
                        tripleLookup rho x
                           | Just w <- lookup rho x
                           , Just y <- lookup rho w
                           , Just z <- lookup rho y
                           = z
                        tripleLookup _ x = handleFailure x
        \end{verbatim}
        \end{center}    
    \caption{Pattern guards swoop in to restore the Nice~Properties, and all is
    right again.} 
    \Description{An implementation of nameOf using pattern guards.}
    \label{fig:guardtriplelookup}
    \end{figure}

    Guards appear as a comma-separated list between the~\tt{|} and the~\tt{=}.
    Each guard has a pattern, followed by~\tt{<-}, then an expression. The guard
    is evaluated by evaluating the expression and testing if the pattern matches
    with the result. If it does, the next guard is evaluated in an
    environment that includes the bindings introduced by evaluating guards
    before it. If the match fails, the program skips evaluating the rest of the
    branch and falls through to the next one. As a bonus, a guard can simply be
    a Boolean expression which the program tests the same way it would a side
    condition, so guards subsume side conditions! 
    
    If you need further convincing on why programmers want for guards, look no
    further than Erwig~\& Peyton Jones's~\it{Pattern Guards and Transformational
    Patterns}~\citep{guardproposal}, the proposal for pattern guards in GHC: the
    authors show several other examples where guards drastically simplify
    otherwise-maddening code. 
    
    Pattern guards enable expressions within guards to utilize names bound in
    preceding guards, enabling imperative pattern-matched steps with expressive
    capabilities akin to Haskell's~\tt{do} notation. It should come as no
    surprise that pattern guards are built in to GHC.~\raggedbottom
\subsubsection{Or-patterns}
    I~conclude our tour of extensions to pattern matching with or-patterns,
    which are built in to OCaml. Let's consider a final example. I~have a type
    \tt{token} which represents an item or location in a video game and how much
    fun it is, and I~need to quickly know what game it's from and how much fun
    I'd have playing it. To do so, I'm going to write a function
    \tt{game\_of\_token} in OCaml. HThe~\tt{token} type and the type and
    algebraic laws for~\tt{game\_of\_token} are in Figure~\ref{fig:gotlaws}.
  
% \begin{minipage}[t]{\textwidth}
\begin{figure}[t]
    \centering 
    \begin{footnotesize}
      \begin{verbatim}  
type funlevel = int

type token = BattlePass  of funlevel | ChugJug    of funlevel | TomatoTown of funlevel
           | HuntersMark of funlevel | SawCleaver of funlevel
           | MoghLordOfBlood of funlevel | PreatorRykard of funlevel
           ... other tokens ...
                    

game_of_token : token -> string * funlevel

game_of_token t == ("Fortnite", f),  where t is any of 
                                      BattlePass f, 
                                      ChugJug f, or
                                      TomatoTown f
game_of_token t == ("Bloodborne", 2 * f), 
                                      where t is any of 
                                        HuntersMark f or 
                                        SawCleaver f
game_of_token t == ("Elden Ring", 3 * f), 
                                      where t is any of 
                                        MoghLordOfBlood f or  
                                        PreatorRykard f
game_of_token _ == ("Irrelevant", 0), otherwise
      \end{verbatim}
    \end{footnotesize}
    \Description{Type and laws for game_of_token.} 
  % \end{minipage}  
    \caption{Type and laws for~\tt{game\_of\_token}, which make helpful use of
    "where."}
    % \ref{fig:gotlaws}
\label{fig:gotlaws}
\end{figure}        
        
        I~can write code for~\tt{game\_of\_token} in OCaml using bare patterns
        (Figure~\ref{fig:baregot}), but I'm dissatisfied with how it fails the
        \hyperref[p1]{Lawlike} and~\hyperref[p2]{Single copy} Nice~Properties:
        it is visually different from the algebraic laws, and it has many
        duplicated right-hand sides.
        
        \begin{figure}
            \begin{center}
                \begin{verbatim}
        let game_of_token token = match token with 
            | BattlePass f      -> ("Fortnite", f)
            | ChugJug f         -> ("Fortnite", f)
            | TomatoTown f      -> ("Fortnite", f)
            | HuntersMark f     -> ("Bloodborne", 2 * f)
            | SawCleaver f      -> ("Bloodborne", 2 * f)
            | MoghLordOfBlood f -> ("Elden Ring", 3 * f)
            | PreatorRykard f   -> ("Elden Ring", 3 * f)
            | _                 -> ("Irrelevant", 0)
                \end{verbatim}
            \end{center}    

        \caption{\tt{game\_of\_token}, with redundant right-hand sides,
        should raise a red flag.} 
        \Description{An implementation of game_of_token using bare patterns.}
        \label{fig:baregot}
        \end{figure}

        I~could try to use a couple of helper functions to reduce clutter, so I
        do so and come up with the code in Figure~\ref{fig:helpergot}. It looks
        ok, but I'm still hurting for Nice~Property~\hyperref[p2]{Lawlike}, and
        now I've lost the~\hyperref[p3]{Call-free} Property, as well.

        \begin{figure}
            \begin{center}
                \begin{verbatim}
                  let fortnite   f = ... complicated   ... in
                  let bloodborne f = ... complicated'  ... in
                  let eldenring  f = ... complicated'' ... in
                  match token with
                  | BattlePass f -> fortnite f
                  ... and so on ...                
                \end{verbatim}
            \end{center}    
        \caption{\tt{game\_of\_token} with helpers is somewhat better, but I'm
        not satisfied with it.} 
        \Description{An implementation of
        game_of_token using helper functions.}
        \label{fig:helpergot}
        \end{figure}

        Once again, an extension comes to the rescue.~\it{Or-patterns} condense
        multiple patterns that share a right-hand side, and when any one of the
        patterns matches with the scrutinee, the right-hand side is evaluated
        with the bindings created by that pattern. I~exploit or-patterns in
        Figure~\ref{fig:orgot} to restore the Nice~Properties and eliminate much
        of the uninteresting code that appeared in
        \ref{fig:baregot}~and~\ref{fig:helpergot}. 

    \begin{figure}
    \begin{center}
    \begin{verbatim}
let game_of_token token = match token with 
    | BattlePass f | ChugJug f | TomatoTown f  -> ("Fortnite", f)
    | HuntersMark f | SawCleaver f             -> ("Bloodborne", 2 * f)
    | MoghLordOfBlood f | PreatorRykard f      -> ("Elden Ring", 3 * f)
    | _                                        -> ("Irrelevant", 0)
    \end{verbatim}
    \end{center}    
    \caption{Or-patterns condense~\tt{game\_of\_token}
    significantly, and it is easier to read line-by-line.} 
    \Description{An
    implementation of game_of_token using or-patterns.}
    \label{fig:orgot}
    \end{figure}

    In addition to the inherent appeal of brevity, or-patterns serve to
    concentrate complexity at a single juncture and create single points of
    truth, restoring the~\hyperref[p1]{Lawlike} and~\hyperref[p2]{Single copy}
    properties. 
    
    %%  And by eliminating helper functions, like the ones in
    %%  Figure~\ref{fig:helpergot}, or-patterns actually boost performance.
    %%
    %%  Not with any self-respecting compiler.  ---NR
      
    \subsubsection{Wrapping up pattern matching and extensions}
    
    I~have presented three popular extensions that make pattern matching more
    expressive and how to use them effectively. Earlier, though, you might have
    noticed a problem. Say I~face a decision-making problem that obliges me to
    use~\emph{all} of these extensions. When picking a language to do so, I am
    stuck! No major functional language has all three of these extensions.
    Remember when I~had to switch from OCaml to Haskell to use guards, and back
    to OCaml for or-patterns? The two extensions are mutually exclusive in
    Haskell and OCaml, and also Scala, Erlang/Elixir, Rust, F\#, and
    Agda~\citep{haskell, ocaml, scala, erlang, elixir, rust, fsharp, agda}. 


    I~find the extension story somewhat unsatisfying. At the very least, I~want
    to be able to use pattern matching, with the extensions I~want, in a single
    language. Or, I~want an alternative that gives me the expressive power of
    pattern matching with these extensions. 

\end{document}
